\section{Complete tables for teaching, libraries, energy and quality}
These four tables reuse existing public interfaces with different layouts.
All values are illustrative; scales, units and missing observations remain explicit.
Standalone files are in \key{docs/examples/}; only the final case needs the optional distribution library.

\subsection{Course projects: completion and status}
Completed and total task counts are listed separately. Percentages are calculated before input, and every bar shares 0--100\%. Status colors are assigned explicitly.

\begin{demo}{Course projects: completion and status}
\begingroup
\gradbarssetup{min=0,max=100,width=42mm,height=6pt,
  precision=0,unit={\%},label=outside,label width=3em,
  label gap=4pt,baseline=text,table legend=true}
\renewcommand{\arraystretch}{1.5}
\begin{center}
\begin{tabular}{lrl}
\toprule
\rowcolor{gradbarsTrack}
Course & Tasks & \gradtablelegendhead[
  stack colors={teal,orange,blue,gray}]{Completion}
  {Done,Active,Planned,Unmeasured}\\
\midrule
Geometry & 18/25 & \gradbar[color=orange]{72}\\
Data & 24/25 & \gradbar[color=orange]{96}\\
Methods & 25/25 & \gradbar[color=teal]{100}\\
Design & 0/25 & \gradbar[color=blue]{0}\\
Seminar & -- & \gradbar[color=gray]{NA}\\
\bottomrule
\end{tabular}
\end{center}
\endgroup
\end{demo}

A real zero retains its 0\% label; an unmeasured record stays missing. The literal percent unit adds a suffix and does not divide task counts.
\par\smallskip\noindent Standalone source: \key{docs/examples/course-progress.tex}.

\clearpage

\subsection{Library loans: monochrome channel shares}
Named categories share textures between the legend and stacked segments. Desk, locker and delivery are mutually exclusive channels in this example.

\begin{demo}{Library loans: monochrome channel shares}
\begingroup
\gradbarscategory{Desk}{blue}{diagonal}
\gradbarscategory{Locker}{teal}{dots}
\gradbarscategory{Delivery}{orange}{crosshatch}
\gradbarssetup{print,min=0,max=300,width=70mm,height=14pt,
  stack names={Desk,Locker,Delivery},segment labels=percent,
  precision=0,label=none,table legend=true}
\renewcommand{\arraystretch}{1.8}
\begin{center}
\begin{tabular}{@{}lrl@{}}
\toprule
Branch & Loans & \gradtablelegendhead[print]
  {Fulfilment share}{Desk,Locker,Delivery}\\
\midrule
North & 300 & \gradstack{180,90,30}\\
South & 300 & \gradstack{120,90,90}\\
East & 300 & \gradstack{60,180,60}\\
West & 300 & \gradstack{150,105,45}\\
\bottomrule
\end{tabular}
\end{center}
\endgroup
\end{demo}

Inputs are counts summing to 300 in each row. Segment percentages use that row total. Equal full widths follow from equal totals, not automatic normalization.
\par\smallskip\noindent Standalone source: \key{docs/examples/borrowing-channels.tex}.

\clearpage

\subsection{Building retrofit: grouped headers and energy use}
A grouped header places before/after values under one unit. All dumbbells share the 0--200 scale rather than enlarging each row independently.

\begin{demo}{Building retrofit: grouped headers and energy use}
\begingroup
\begin{center}
\begin{gradtable}[width=\linewidth,align=c,header align=c,
  column sep=3pt,row stretch=1.7,label width=9mm]
\gradcolumn[width=22mm]{building}{Building}
\gradcolumn[type=number,width=16mm,
  options={precision=0}]{before}{Before}
\gradcolumn[type=number,width=16mm,
  options={precision=0}]{after}{After}
\gradcolumn[type=dumbbell,options={min=0,max=200,
  precision=0,compare label=delta,track color=black!35,
  stem width=.9pt}]{change}{Change}
\gradheader{\gradspan{1}{Site}
  \gradspan{2}{Annual use (kWh/m$^2$)}
  \gradspan{1}{Common scale: 0--200}}
\gradrow{Atrium,160,120,160/120}
\gradrow{Studio,130,90,130/90}
\gradrow{Archive,80,76,80/76}
\gradrow{Annex,110,128,110/128}
\end{gradtable}
\end{center}
\endgroup
\end{demo}

Hollow points are before and filled points are after. Delta is after minus before, in kWh/m$^2$. A negative value means lower use in these illustrative observations, not a causal estimate.
\par\smallskip\noindent Standalone source: \key{docs/examples/energy-retrofit.tex}.

\clearpage

\subsection{Production batches: box plots on a shared scale}
Load \key{\usegradbarslibrary{distribution}} in the preamble. Every batch uses 9.80--10.20 mm so location, spread and outliers remain comparable.

\begin{demo}{Production batches: box plots on a shared scale}
\begingroup
\gradbarsstyle{batch-demo}{box range=fixed,min=9.8,max=10.2,
  width=70mm,height=5ex,precision=2,label=outside,
  unit={\,mm},palette=accessible}
\newcommand{\batchplot}[1]{\raisebox{\dimexpr.5ex-.5\height+.5\depth\relax}{%
  \gradboxplot[style=batch-demo]{#1}}}
\renewcommand{\arraystretch}{2.2}
\begin{center}
\begin{tabular}{@{}lrp{94mm}@{}}
\toprule
Batch & Samples & 9.80--10.20 mm / median\\
\midrule
A & 9 & \batchplot{9.96,9.98,9.99,10,10.01,10.02,10.03,10.04,10.06}\\
B & 9 & \batchplot{9.85,9.95,9.98,10,10.01,10.02,10.04,10.08,10.17}\\
C & 9 & \batchplot{10.03,10.04,10.05,10.05,10.06,10.07,10.07,10.08,10.09}\\
Pending & -- & \batchplot{NA}\\
\bottomrule
\end{tabular}
\end{center}
\endgroup
\end{demo}

Right-hand labels show medians. Whiskers end at observed values inside the 1.5 IQR fences, and hollow points retain outliers. The display range is not an acceptance interval; whiskers are not confidence intervals.
\par\smallskip\noindent Standalone source: \key{docs/examples/batch-dimensions.tex}.
